\chapter{Crypto Medium-Frequency Strategies}\label{ch:s1:crypto-medium-frequency-strategies}

In 2021 the funding on Binance's Bitcoin perpetual averaged 30.6\% a year, paid by longs to shorts: a book long spot bitcoin and short the perpetual,
with no view on the price, collected it. In November 2022 FTX.com, one of the largest venues on which such books were run, stopped paying withdrawals;
at the petition its debtors located assets worth about a fifth of what customers were owed. Crypto's medium-frequency strategies are carry, momentum
and flows, like other markets', with one addition: the venue itself can fail. On Binance's data the carry book earned 25.3\% on capital in 2021 and
3.3\% in 2022; a book spread over four venues, each failing with a 5\% chance a year and losing what FTX lost, gives up 20.2\% of its capital when one
fails. The build is \texttt{firm.cryptomf}.

\section{Funding and basis carry}

\begin{definition}[Funding carry, basis carry]\label{def:s1:crypto-medium-frequency-strategies:carry}
\emph{Funding carry}\index{funding carry} is the return of a position long spot and short a perpetual future of the same coin, which receives the
perpetual's funding payments when they are positive (longs pay shorts) and pays them when negative, with no net exposure to the price. \emph{Basis
carry}\index{basis carry} is the same trade with a dated future: long spot, short the future, earning the future's premium over spot as it converges
at expiry.
\end{definition}

\begin{dated}{2026-09}{Perpetual funding on Binance}\label{dat:s1:crypto-medium-frequency-strategies:funding}
Binance's funding rate for a perpetual is the average premium of the perpetual over its index plus a clamped interest term, paid every eight hours at
00:00, 08:00 and 16:00 UTC; for most contracts the interest term is 0.01\% per interval. Book~3, chapter~17 describes the mechanism and its premium
index.
\end{dated}

The carry book's return on capital (\cref{lst:s1:crypto-medium-frequency-strategies:data}) is the year's funding, less trading costs (0.2\% a year,
assumed), divided by the capital it ties up: the spot position plus the perpetual's margin (20\% of notional, assumed).

\begin{center}
\small
\begin{tabular}{@{}lccccccc@{}}
\toprule
year & 2020 & 2021 & 2022 & 2023 & 2024 & 2025 & 2026\\
\midrule
funding, mean annualised & 17.19\% & 30.61\% & 4.16\% & 7.87\% & 11.92\% & 5.13\% & 2.94\%\\
share of negative rates & 14.3\% & 7.3\% & 22.1\% & 10.1\% & 8.4\% & 12.9\% & 26.1\%\\
carry book, return on capital & 14.2\% & 25.3\% & 3.3\% & 6.4\% & 9.8\% & 4.1\% & 2.3\%\\
\bottomrule
\end{tabular}
\end{center}

2026 runs to 24 September. The carry follows the market's demand for leverage (\cref{fig:s1:crypto-medium-frequency-strategies:funding}): high in the
rising market of 2021, small after 2022's crash, recovering in 2024. Schmeling, Schrimpf and Todorov found the same in dated futures' basis, sometimes
above 40\% a year, and traced it to smaller, trend-chasing investors' demand for leverage meeting arbitrage capital limited by regulatory and margin
frictions. The carry is the price of that scarcity. In 2024 it was close to the interest term of 0.01\% per interval ($0.0001 \times 3 \times 365 = 11\%$ a
year); since then it has been below it.

\begin{omfigure}
\begin{tikzpicture}
\begin{axis}[width=10.5cm,height=5cm,ybar,bar width=10pt,xlabel={\small year},ylabel={\small \% a year},tick label style={font=\footnotesize},
  symbolic x coords={2020,2021,2022,2023,2024,2025,2026},xtick=data,ymin=0,ymax=35,enlarge x limits=0.08,grid=major,grid style={gray!25},
  legend columns=2,legend style={font=\scriptsize,at={(0.5,-0.3)},anchor=north,draw=none,fill=none,
  /tikz/every even column/.append style={column sep=8pt}},area legend,table/col sep=comma]
\addplot[fill=omExo!50,draw=omExo] table[x=year,y=funding_pct]{figdata/strategies-1/27-crypto-medium-frequency-strategies/funding.csv};
\addplot[fill=omThm!60,draw=omThm] table[x=year,y=carry_pct]{figdata/strategies-1/27-crypto-medium-frequency-strategies/funding.csv};
\legend{mean funding (annualised),carry book's return on capital}
\end{axis}
\end{tikzpicture}

\omcaption{Binance BTCUSDT perpetual: the mean funding rate by calendar year, annualised, and the return on capital of a book long spot and short the
perpetual (20\% margin, 0.2\% a year of costs, assumed); 2026 to 24 September. Data: Binance's funding-rate endpoint via Book~3;
\texttt{s1\_crypto.funding}.}\label{fig:s1:crypto-medium-frequency-strategies:funding}
\end{omfigure}

\section{Cross-sectional momentum in coins}

Liu and Tsyvinski found that cryptocurrency returns are not exposed to the usual stock-market and macroeconomic factors but are predicted by factors
specific to crypto: a strong \omterm{def:s1:trend-following:trend}{time-series momentum} effect and proxies for investor attention. Liu, Tsyvinski and Wu found that three factors, the crypto
market, size and momentum, account for the cross-section of expected coin returns. \texttt{firm.cryptomf} simulates six years of 50 coins with a 60\%
market factor, 80\% specific volatility with fat tails (a Student-$t$ with three degrees of freedom), and a planted persistent drift with a half-life of 60
days, and runs a weekly long--short book on past returns, top fifth against bottom fifth:

\begin{center}
\small
\begin{tabular}{@{}lccc@{}}
\toprule
lookback & 7 days & 28 days & 56 days\\
\midrule
Sharpe ratio before costs & 0.85 & 1.56 & 1.29\\
after 10 bp per unit traded & 0.42 & 1.35 & 1.13\\
after 20 bp per unit traded & $-0.01$ & 1.13 & 0.97\\
\bottomrule
\end{tabular}
\end{center}

A week of history is mostly noise at these volatilities; four weeks measure the drift better and trade less. Costs in crypto vary widely by venue and
coin, and the small coins where momentum is strongest are the most expensive to trade.

\section{On-chain data signals}

\begin{definition}[On-chain signal]\label{def:s1:crypto-medium-frequency-strategies:onchain}
An \emph{on-chain signal}\index{on-chain signal} is a trading signal built from a blockchain's public records: flows of coins to and from exchanges'
known addresses, balances of large holders, activity of addresses, issuance and burning, available to anyone who reads the chain (Book~3,
chapter~26).
\end{definition}

The classic hypothesis: coins sent to exchanges are about to be sold. The chapter plants it: each coin's net inflow, standardised, has a correlation of
$-0.03$ with the next day's specific return. Measured as a rank IC across the 50 coins, the z-scored inflow (against its past 30 days) gives $-0.031$
($t = -10.2$) over six years. The planted effect is recovered. In real data the mapping of addresses to exchanges is incomplete and changes, and flows
between an exchange's own wallets look like deposits; the scaffolding in \texttt{firm.cryptomf} is where such a signal is tested, not evidence that one
works.

\section{Venue and custody risk}

\begin{definition}[Venue risk]\label{def:s1:crypto-medium-frequency-strategies:venue}
\emph{Venue risk}\index{venue risk} is the risk of losing assets held at a trading venue or custodian through its failure, fraud, hacking or freezing of
withdrawals; in crypto it is borne by the customer, because exchanges often hold customers' assets themselves.
\end{definition}

FTX.com's debtors' analysis gives the size of such a loss. At the petition, customer payables summed to about \$11.2 billion and the located assets to
about \$2.2 billion: a shortfall of 80.8\%. In the eleven days of November 2022 before the petition, customers withdrew \$7.0 billion more than they
deposited, \$3.3 billion net on 7 November alone. Book~3, chapter~15 tells the story. For a carry book the lesson is arithmetic
(\cref{lst:s1:crypto-medium-frequency-strategies:venues}): with venues that fail independently with probability $p$ a year and lose a share $L$ of the
assets on them, the expected loss is $pL$ whatever the number of venues, and spreading assets only changes its distribution.

\begin{center}
\small
\begin{tabular}{@{}lcccc@{}}
\toprule
venues (equal shares), $p = 5\%$, $L = 80.8\%$ & 1 & 2 & 4 & 8\\
\midrule
loss when one venue fails (share of capital) & 80.8\% & 40.4\% & 20.2\% & 10.1\%\\
chance of at least one failure in a year & 5.0\% & 9.9\% & 18.7\% & 34.0\%\\
expected loss a year & 4.1\% & 4.1\% & 4.1\% & 4.1\%\\
99th percentile of the yearly loss & 80.8\% & 40.4\% & 40.4\% & 20.2\%\\
\bottomrule
\end{tabular}
\end{center}

Against the carry, an expected loss of about 4\% a year is decisive: 2021's 25.3\% would net 21.3\%, 2022's 3.3\% would be a loss
(\cref{fig:s1:crypto-medium-frequency-strategies:venues}). Diversification across venues makes the bad year survivable; it does not make the strategy
better. What lowers the expected loss is the choice of venues and custody: segregated custody, assets held off the exchange until needed, and the
speed to withdraw when a venue weakens.

\begin{omfigure}
\begin{tikzpicture}
\begin{axis}[width=10cm,height=5cm,ybar,bar width=12pt,xlabel={\small number of venues},ylabel={\small \% of capital, or chance},
  tick label style={font=\footnotesize},symbolic x coords={1,2,4,8},xtick=data,ymin=0,ymax=90,enlarge x limits=0.2,grid=major,
  grid style={gray!25},legend columns=3,legend style={font=\scriptsize,at={(0.5,-0.3)},anchor=north,draw=none,fill=none,
  /tikz/every even column/.append style={column sep=8pt}},area legend,table/col sep=comma]
\addplot[fill=omThm!60,draw=omThm] table[x=k,y=one_pct]{figdata/strategies-1/27-crypto-medium-frequency-strategies/venues.csv};
\addplot[fill=omDef!60,draw=omDef] table[x=k,y=any_pct]{figdata/strategies-1/27-crypto-medium-frequency-strategies/venues.csv};
\addplot[fill=omExo!50,draw=omExo] table[x=k,y=mean_pct]{figdata/strategies-1/27-crypto-medium-frequency-strategies/venues.csv};
\legend{loss when one fails,chance of a failure a year,expected loss a year}
\end{axis}
\end{tikzpicture}

\omcaption{A carry book's capital spread equally over one to eight venues, each failing independently with a 5\% chance a year (assumed) and losing
80.8\% of the assets on it (the FTX.com petition-time shortfall): the loss when one venue fails, the chance of at least one failure in a year, and the
expected yearly loss. Data: \texttt{s1\_crypto.venues}.}\label{fig:s1:crypto-medium-frequency-strategies:venues}
\end{omfigure}

\section{Strategy files}

\begin{strategyfile}{Perpetual funding carry}\label{strat:s1:crypto-medium-frequency-strategies:funding}
\sfield{Who pays you, and why} Leveraged longs who pay funding to hold perpetuals.
\sfield{Instruments and venues} Spot coins and their perpetuals on the largest venues.
\sfield{Signal} Positive expected funding; exit when it turns persistently negative.
\sfield{Sizing and execution} Long spot, short perpetual in equal notional; margin kept above maintenance with a buffer.
\sfield{Costs} Trading fees, borrowing if spot is financed, and \omterm{def:s1:crypto-medium-frequency-strategies:venue}{venue risk}.
\sfield{How it dies} Arbitrage capital compresses funding; venue failure.
\sfield{Horizon, capacity, infrastructure} Months; capacity limited by open interest; custody and margin monitoring.
\sfield{Backtest honestly} Funding as paid, interval by interval; fees; the capital actually tied up; a venue-loss charge.
\sfield{Sources} Binance funding data (Book~3); this chapter: 25.3\% on capital in 2021, 3.3\% in 2022, before \omterm{def:s1:crypto-medium-frequency-strategies:venue}{venue risk}.
\end{strategyfile}

\begin{strategyfile}{Dated-futures basis carry}\label{strat:s1:crypto-medium-frequency-strategies:basis}
\sfield{Who pays you, and why} Investors wanting leveraged long exposure through futures.
\sfield{Instruments and venues} Dated futures on regulated and crypto venues, against spot.
\sfield{Signal} The annualised basis to expiry.
\sfield{Sizing and execution} Long spot, short the future, held to expiry.
\sfield{Costs} Fees; capital for margin; the roll into the next expiry.
\sfield{How it dies} More arbitrage capital; regulation that eases it.
\sfield{Horizon, capacity, infrastructure} Weeks to months.
\sfield{Backtest honestly} Basis at executable prices; margin at the venue's rules.
\sfield{Sources} Schmeling, Schrimpf and Todorov: carry sometimes above 40\% a year.
\end{strategyfile}

\begin{strategyfile}{Coin cross-sectional momentum}\label{strat:s1:crypto-medium-frequency-strategies:momentum}
\sfield{Who pays you, and why} Investors who underreact, then chase; attention flows.
\sfield{Instruments and venues} Spot or perpetuals on liquid coins.
\sfield{Signal} Past three- to eight-week returns, ranked.
\sfield{Sizing and execution} Weekly long--short quintiles, volatility-scaled.
\sfield{Costs} Large for small coins.
\sfield{How it dies} Crashes and reversals; delistings.
\sfield{Horizon, capacity, infrastructure} Weeks; clean price histories including dead coins.
\sfield{Backtest honestly} Survivorship: include delisted coins; costs by coin.
\sfield{Sources} Liu and Tsyvinski (2021); Liu, Tsyvinski and Wu (2022); this chapter: 1.35 after 10 bp at four weeks, synthetic.
\end{strategyfile}

\begin{strategyfile}{On-chain flow signal}\label{strat:s1:crypto-medium-frequency-strategies:onchain}
\sfield{Who pays you, and why} Holders about to sell, whose deposits to exchanges are public before their orders.
\sfield{Instruments and venues} Coins with well-labelled exchange addresses.
\sfield{Signal} Z-scored net inflows to exchanges.
\sfield{Sizing and execution} Short-horizon tilts; small.
\sfield{Costs} Data labelling; turnover.
\sfield{How it dies} Internal transfers mistaken for deposits; wider use of the same data.
\sfield{Horizon, capacity, infrastructure} Days; a node or a data provider and address labels.
\sfield{Backtest honestly} Labels as they were known at the time; block times, not later-indexed times.
\sfield{Sources} No performance figure verified; the chapter's synthetic test.
\end{strategyfile}

\begin{strategyfile}{Funding-rate mean reversion}\label{strat:s1:crypto-medium-frequency-strategies:meanrev}
\sfield{Who pays you, and why} Crowded leveraged positioning that unwinds.
\sfield{Instruments and venues} Perpetuals.
\sfield{Signal} Funding far above or below its recent range.
\sfield{Sizing and execution} Fade extreme funding, sized small.
\sfield{Costs} Funding paid while waiting; liquidation cascades.
\sfield{How it dies} Trends that persist while funding stays extreme.
\sfield{Horizon, capacity, infrastructure} Days.
\sfield{Backtest honestly} Funding interval by interval; liquidation risk at the venue's margin rules.
\sfield{Sources} No performance figure verified.
\end{strategyfile}

\begin{strategyfile}{Venue-diversified carry}\label{strat:s1:crypto-medium-frequency-strategies:venue}
\sfield{Who pays you, and why} As for \omterm{def:s1:crypto-medium-frequency-strategies:carry}{funding carry}.
\sfield{Instruments and venues} The same trade on several venues.
\sfield{Signal} Funding by venue, net of each venue's risk.
\sfield{Sizing and execution} Capital spread by venue quality; withdrawal limits and triggers.
\sfield{Costs} Lower netting, more capital idle.
\sfield{How it dies} Correlated failures when one venue's collapse spreads.
\sfield{Horizon, capacity, infrastructure} Months; treasury operations across venues.
\sfield{Backtest honestly} A venue-failure charge; failures in the sample.
\sfield{Sources} FTX debtors' analysis (Doc 792-1); this chapter: 20.2\% of capital lost when one of four venues fails.
\end{strategyfile}

\section{Tutorial: until the exchange failed}\label{tut:s1:crypto-medium-frequency-strategies:tut}

\begin{tutorial}
\textbf{Goal.} Turn Binance's funding statistics into a carry book's yearly returns, measure FTX's loss given failure from its debtors' analysis,
simulate coin momentum and exchange flows, and size a carry book across venues. \textbf{End state:} the tables and the two figures.
\begin{tutsteps}
\item \textbf{The real data}: funding by year and the FTX balances and flows.
\omcode{code/strategies-1/27-crypto-medium-frequency-strategies/python/s1_crypto.py}{38}{54}{Funding carry by year and the FTX shortfall.}\label{lst:s1:crypto-medium-frequency-strategies:data}
\item \textbf{Momentum, flows and venues}.
\omcode{code/firm/cryptomf/firm_cryptomf.py}{65}{100}{Coin momentum, flow IC and venue losses.}\label{lst:s1:crypto-medium-frequency-strategies:venues}
\item \textbf{Run} \texttt{funding()}, \texttt{ftx()}, \texttt{momentum\_table()}, \texttt{flows()} and \texttt{venues()}, and \texttt{fig\_crypto.py}.
\end{tutsteps}
\textbf{What to change next.} Correlate venue failures (a failure raises the others' probability); fetch funding interval by interval and compute
drawdowns; add a delisting process to the coin universe and see momentum's survivorship bias.
\end{tutorial}

\section{Build: crypto strategies}\label{bld:s1:crypto-medium-frequency-strategies:cryptomf}

\begin{build}
\bfield{Purpose} Carry accounting, a coin universe with momentum and flows, and venue-failure losses.
\bfield{Interface} \texttt{carry\_return(funding, margin, cost)}, \texttt{basis\_carry(fut, spot, days)}, \texttt{CoinConfig(\dots)},
\texttt{simulate\_coins(cfg, rng)}, \texttt{momentum\_book(r, lookback, hold, q, cost)}, \texttt{flow\_ic(flow, r, window)}, \texttt{venue\_losses(k,
p, lgd, years, rng)}.
\bfield{Rules} Carry on the capital tied up; signals from past data; venue failures independent unless stated.
\bfield{Acceptance tests} \texttt{code/firm/cryptomf/tests/}: carry and basis by hand, momentum and the flow IC on a planted universe, the mean and support
of venue losses.
\bfield{Stretch} Correlated failures; interval-level funding; delistings.
\end{build}

\begin{omsources}
\item Y.~Liu and A.~Tsyvinski, ``Risks and returns of cryptocurrency'', \emph{Review of Financial Studies} 34(6), 2021.
\item Y.~Liu, A.~Tsyvinski and X.~Wu, ``Common risk factors in cryptocurrency'', \emph{Journal of Finance} 77(2), 2022.
\item M.~Schmeling, A.~Schrimpf and K.~Todorov, ``Crypto carry'', \emph{Management Science}, 2026.
\item FTX Debtors, \emph{Preliminary Analysis of Shortfalls}, Doc 792-1, In re FTX Trading Ltd., 22-11068 (Bankr. D. Del.), March 2023.
\item Binance, BTCUSDT perpetual funding rates (public endpoint) and funding-rate documentation.
\end{omsources}

\section{Exercises}

\begin{exercise}[$\star$]\label{exo:s1:crypto-medium-frequency-strategies:1}
A funding rate of 0.01\% per eight-hour interval, paid three times a day: what is it annualised? And 0.1\%?
\end{exercise}

\begin{exercise}[$\star$]\label{exo:s1:crypto-medium-frequency-strategies:2}
Funding averages 30.61\% a year; costs are 0.2\% a year; the perpetual's margin is 20\% of notional. What does the book earn on its capital?
\end{exercise}

\begin{exercise}[$\star$]\label{exo:s1:crypto-medium-frequency-strategies:3}
Why is a book long spot and short the perpetual not exposed to the coin's price, and what is it exposed to?
\end{exercise}

\begin{exercise}[$\star\star$]\label{exo:s1:crypto-medium-frequency-strategies:4}
Venues fail independently with a 5\% chance a year and lose 80.8\% of the assets on them. What is the expected yearly loss, and the chance of at least
one failure among four venues? Among eight?
\end{exercise}

\begin{exercise}[$\star\star$]\label{exo:s1:crypto-medium-frequency-strategies:5}
Net of an expected venue loss of 4.0\% a year, what did the carry book earn in 2021 and in 2022?
\end{exercise}

\begin{exercise}[$\star\star$]\label{exo:s1:crypto-medium-frequency-strategies:6}
Why does the one-week momentum book lose after 20 basis points while the four-week book does not?
\end{exercise}

\begin{exercise}[$\star\star\star$]\label{exo:s1:crypto-medium-frequency-strategies:7}
\emph{Coding.} Run \texttt{venues\_correlated()}: in a year with a failure, each other venue also fails with probability 30\%. Compare the expected
loss and the 99th percentile with the independent case for four venues.
\end{exercise}

\begin{exercise}[$\star\star\star$]\label{exo:s1:crypto-medium-frequency-strategies:8}
\emph{Find the flaw.} ``Our \omterm{def:s1:crypto-medium-frequency-strategies:carry}{funding carry} earned 18\% a year for five years with a maximum drawdown of 2\%: it is market-neutral and nearly riskless.''
\end{exercise}

\section{Problem: Until the Exchange Failed}

\begin{problem}[{Weekend problem --- carry, momentum and the venue}]\label{pb:s1:crypto-medium-frequency-strategies:1}
Binance's funding data, the FTX debtors' analysis, the synthetic coins and the public record.

\medskip\noindent\textbf{Part I --- Carry.}
\begin{enumerate}
\item Define \omterm{def:s1:crypto-medium-frequency-strategies:carry}{funding carry} and \omterm{def:s1:crypto-medium-frequency-strategies:carry}{basis carry}.
\item How is Binance's funding rate set, and when is it paid?
\item Give the carry book's return on capital by year.
\item What did Schmeling, Schrimpf and Todorov find about crypto carry?
\end{enumerate}

\medskip\noindent\textbf{Part II --- Momentum and flows.}
\begin{enumerate}[resume]
\item What did Liu and Tsyvinski, and Liu, Tsyvinski and Wu find?
\item Describe the synthetic coin universe.
\item Give the momentum book's Sharpe ratios by lookback and cost.
\item Define an \omterm{def:s1:crypto-medium-frequency-strategies:onchain}{on-chain signal}; what did the flow test find, and what does it not show?
\end{enumerate}

\medskip\noindent\textbf{Part III --- Venues.}
\begin{enumerate}[resume]
\item Define \omterm{def:s1:crypto-medium-frequency-strategies:venue}{venue risk}.
\item What do FTX.com's petition-time balances and November flows show?
\item Give the loss when one of $k$ venues fails and the chance of a failure a year.
\item Why does diversification not change the expected loss?
\end{enumerate}

\medskip\noindent\textbf{Part IV --- The verdict.}
\begin{enumerate}[resume]
\item State the \emph{named result}: the \omterm{def:s1:crypto-medium-frequency-strategies:carry}{funding carry}'s yearly return and the venue-diversified book's loss when one venue fails.
\item What does lower the expected venue loss?
\item How would you backtest \omterm{def:s1:crypto-medium-frequency-strategies:carry}{funding carry} honestly?
\item Which strategy file carries the most \omterm{def:s1:crypto-medium-frequency-strategies:venue}{venue risk}?
\item How does this chapter's carry relate to chapter~20's?
\item What did the carry pay for in 2021?
\item Why did funding fall after 2022?
\item In one sentence: what is a crypto carry trader paid for?
\end{enumerate}
\end{problem}

\section{Interview questions}

\begin{interviewq}[$\star$ \iqroles{researcher}]\label{iq:s1:crypto-medium-frequency-strategies:1}
What is a perpetual future's funding rate, and why does it exist?
\end{interviewq}

\begin{interviewq}[$\star\star$ \iqroles{trader}]\label{iq:s1:crypto-medium-frequency-strategies:2}
Funding on a coin is 60\% a year. Walk through the trade and its risks.
\end{interviewq}

\begin{interviewq}[$\star\star$ \iqroles{risk}]\label{iq:s1:crypto-medium-frequency-strategies:3}
How would you set limits on the assets a firm keeps on each crypto venue?
\end{interviewq}

\begin{interviewq}[$\star\star$ \iqroles{developer}]\label{iq:s1:crypto-medium-frequency-strategies:4}
What does a system running a carry book across several venues need to monitor?
\end{interviewq}

\begin{interviewq}[$\star\star$ \iqroles{researcher}]\label{iq:s1:crypto-medium-frequency-strategies:5}
How would you test coin momentum without survivorship bias?
\end{interviewq}

\begin{interviewq}[$\star\star\star$ \iqroles{researcher}]\label{iq:s1:crypto-medium-frequency-strategies:6}
A book spreads its capital over $k$ venues, each failing independently with probability $p$ a year and losing a share $L$. Derive the mean and variance
of the yearly loss, and the probability that it is at least $mL/k$ for an integer $m$.
\end{interviewq}
